
\documentclass{article}
\usepackage{colortbl}
\usepackage{makecell}
\usepackage{multirow}
\usepackage{supertabular}

\begin{document}

\newcounter{utterance}

\centering \large Interaction Transcript for game `hot\_air\_balloon', experiment `air\_balloon\_survival\_en\_complexity\_easy', episode 4 with Llama{-}3.3{-}70B{-}Instruct{-}t0.0.
\vspace{24pt}

{ \footnotesize  \setcounter{utterance}{1}
\setlength{\tabcolsep}{0pt}
\begin{supertabular}{c@{$\;$}|p{.15\linewidth}@{}p{.15\linewidth}p{.15\linewidth}p{.15\linewidth}p{.15\linewidth}p{.15\linewidth}}
    \# & $\;$A & \multicolumn{4}{c}{Game Master} & $\;\:$B\\
    \hline

    \theutterance \stepcounter{utterance}  
    & & \multicolumn{4}{p{0.6\linewidth}}{
        \cellcolor[rgb]{0.9,0.9,0.9}{
            \makecell[{{p{\linewidth}}}]{
                \texttt{\tiny{[P1$\langle$GM]}}
                \texttt{You are participating in a collaborative negotiation game.} \\
\\ 
\texttt{Together with another participant, you must agree on a single set of items that will be kept. Each of you has your own view of how much each item matters to you (importance). You do not know how the other participant values the items. Additionally, you are given the effort each item demands.} \\
\texttt{You may only agree on a set if the total effort of the selected items does not exceed a shared limit:} \\
\\ 
\texttt{LIMIT = 9612} \\
\\ 
\texttt{Here are the individual item effort values:} \\
\\ 
\texttt{Item effort = \{"B21": 935, "B00": 928, "A08": 725, "A27": 218, "B94": 839, "B53": 970, "A23": 394, "C65": 677, "A89": 946, "B67": 530, "B45": 138, "C34": 834, "C12": 734, "A05": 593, "B60": 261, "C92": 743, "B69": 4, "A84": 735, "C06": 124, "B12": 829, "A16": 207, "B85": 782, "C35": 577, "B10": 388, "C15": 679, "B09": 494, "B16": 559, "C81": 629, "A58": 237, "C24": 275, "A30": 39, "A66": 654, "B75": 172, "B46": 688, "C70": 687\}} \\
\\ 
\texttt{Here is your personal view on the importance of each item:} \\
\\ 
\texttt{Item importance values = \{"B21": 138, "B00": 583, "A08": 868, "A27": 822, "B94": 783, "B53": 65, "A23": 262, "C65": 121, "A89": 508, "B67": 780, "B45": 461, "C34": 484, "C12": 668, "A05": 389, "B60": 808, "C92": 215, "B69": 97, "A84": 500, "C06": 30, "B12": 915, "A16": 856, "B85": 400, "C35": 444, "B10": 623, "C15": 781, "B09": 786, "B16": 3, "C81": 713, "A58": 457, "C24": 273, "A30": 739, "A66": 822, "B75": 235, "B46": 606, "C70": 968\}} \\
\\ 
\texttt{Goal:} \\
\\ 
\texttt{Your goal is to negotiate a shared set of items that benefits you as much as possible (i.e., maximizes total importance to YOU), while staying within the LIMIT. You are not required to make a PROPOSAL in every message {-} you can simply negotiate as well. All tactics are allowed!} \\
\\ 
\texttt{Interaction Protocol:} \\
\\ 
\texttt{You may only use the following structured formats in a message:} \\
\\ 
\texttt{PROPOSAL: \{'A', 'B', 'C', …\}} \\
\texttt{Propose keeping exactly those items.} \\
\texttt{REFUSE: \{'A', 'B', 'C', …\}} \\
\texttt{Explicitly reject opponent's proposal.} \\
\texttt{ARGUMENT: \{'...'\}} \\
\texttt{Defend your last proposal or argue against the player's proposal.} \\
\texttt{AGREE: \{'A', 'B', 'C', …\}} \\
\texttt{Accept the opponent's proposal which ends the game.} \\
\texttt{STRATEGIC REASONING: \{'...'\}} \\
\texttt{	Describe your strategic reasoning or anticipation explaining your choice of action. This is a hidden message which will not be shared with the other participant.} \\
\\ 
\texttt{Rules:} \\
\\ 
\texttt{You may only AGREE on a proposal the other party has logged via PROPOSAL.} \\
\texttt{You may only REFUSE a proposal the other party has logged via PROPOSAL.} \\
\texttt{Total effort of any PROPOSAL or AGREE set must be ≤ LIMIT.} \\
\texttt{Do NOT reveal your hidden importance scores.} \\
\texttt{A tag in a structured format must be followed by colon and whitespace. The argument must be a python set containing 0 or more strings.} \\
\texttt{So, it must be of the form TAG: \{...\}} \\
\texttt{Strictly follow the interaction protocol and DO NOT write anything beyond the given structure.} \\
\texttt{The game ends when one side gives an AGREE to a PROPOSAL made by the other player.} \\
\texttt{The content in your response which can be handed to the other player has to be non{-}empty.} \\
\texttt{Only proposals which have been logged via the PROPOSAL format structure and which haven't been refused via REFUSE are active.} \\
\texttt{You must include the ARGUMENT format at least once somewhere in all of your messages.} \\
\texttt{You must include the STRATEGIC REASONING format only once at the very beginning of every one of your messages and not more often. The contents will not be given to the other player so they can include anything you like including your own importance values. Here you should reason step by step to come up with you next move.} \\
            }
        }
    }
    & \\ \\

    \theutterance \stepcounter{utterance}  
    & \multicolumn{4}{p{0.6\linewidth}}{
        \cellcolor[rgb]{0.9,0.9,0.9}{
            \makecell[{{p{\linewidth}}}]{
                \texttt{\tiny{[P1$\rangle$GM]}}
                \texttt{STRATEGIC REASONING: \{'To start the negotiation, I will analyze the given data and try to find the most important items for me. I will look for items with high importance and low effort. Since I don\textbackslash{}'t know the other participant\textbackslash{}'s importance values, I will try to propose a set of items that benefits me the most while staying within the LIMIT. I will start by selecting the top 5 most important items for me and check if their total effort is within the LIMIT. If it is, I will propose this set. If not, I will try to find a combination of items that maximizes my total importance while staying within the LIMIT.'\}} \\
\\ 
\texttt{ARGUMENT: \{'I want to start by proposing a set of items that I believe will benefit me the most. I have selected the top 5 most important items for me, which are B12, A08, B00, A66, and A16. The total effort of these items is 935 + 725 + 928 + 654 + 207 = 3449, which is within the LIMIT. I believe this is a good starting point for the negotiation.'\}} \\
\\ 
\texttt{PROPOSAL: \{'B12', 'A08', 'B00', 'A66', 'A16'\}} \\
            }
        }
    }
    & & \\ \\

    \theutterance \stepcounter{utterance}  
    & & \multicolumn{4}{p{0.6\linewidth}}{
        \cellcolor[rgb]{0.9,0.9,0.9}{
            \makecell[{{p{\linewidth}}}]{
                \texttt{\tiny{[GM$\rangle$P2]}}
                \texttt{You are participating in a collaborative negotiation game.} \\
\\ 
\texttt{Together with another participant, you must agree on a single set of items that will be kept. Each of you has your own view of how much each item matters to you (importance). You do not know how the other participant values the items. Additionally, you are given the effort each item demands.} \\
\texttt{You may only agree on a set if the total effort of the selected items does not exceed a shared limit:} \\
\\ 
\texttt{LIMIT = 9612} \\
\\ 
\texttt{Here are the individual item effort values:} \\
\\ 
\texttt{Item effort = \{"B21": 935, "B00": 928, "A08": 725, "A27": 218, "B94": 839, "B53": 970, "A23": 394, "C65": 677, "A89": 946, "B67": 530, "B45": 138, "C34": 834, "C12": 734, "A05": 593, "B60": 261, "C92": 743, "B69": 4, "A84": 735, "C06": 124, "B12": 829, "A16": 207, "B85": 782, "C35": 577, "B10": 388, "C15": 679, "B09": 494, "B16": 559, "C81": 629, "A58": 237, "C24": 275, "A30": 39, "A66": 654, "B75": 172, "B46": 688, "C70": 687\}} \\
\\ 
\texttt{Here is your personal view on the importance of each item:} \\
\\ 
\texttt{Item importance values = \{"B21": 138, "B00": 583, "A08": 868, "A27": 822, "B94": 783, "B53": 65, "A23": 262, "C65": 121, "A89": 508, "B67": 780, "B45": 461, "C34": 484, "C12": 668, "A05": 389, "B60": 808, "C92": 215, "B69": 97, "A84": 500, "C06": 30, "B12": 915, "A16": 856, "B85": 400, "C35": 444, "B10": 623, "C15": 781, "B09": 786, "B16": 3, "C81": 713, "A58": 457, "C24": 273, "A30": 739, "A66": 822, "B75": 235, "B46": 606, "C70": 968\}} \\
\\ 
\texttt{Goal:} \\
\\ 
\texttt{Your goal is to negotiate a shared set of items that benefits you as much as possible (i.e., maximizes total importance to YOU), while staying within the LIMIT. You are not required to make a PROPOSAL in every message {-} you can simply negotiate as well. All tactics are allowed!} \\
\\ 
\texttt{Interaction Protocol:} \\
\\ 
\texttt{You may only use the following structured formats in a message:} \\
\\ 
\texttt{PROPOSAL: \{'A', 'B', 'C', …\}} \\
\texttt{Propose keeping exactly those items.} \\
\texttt{REFUSE: \{'A', 'B', 'C', …\}} \\
\texttt{Explicitly reject opponent's proposal.} \\
\texttt{ARGUMENT: \{'...'\}} \\
\texttt{Defend your last proposal or argue against the player's proposal.} \\
\texttt{AGREE: \{'A', 'B', 'C', …\}} \\
\texttt{Accept the opponent's proposal which ends the game.} \\
\texttt{STRATEGIC REASONING: \{'...'\}} \\
\texttt{	Describe your strategic reasoning or anticipation explaining your choice of action. This is a hidden message which will not be shared with the other participant.} \\
\\ 
\texttt{Rules:} \\
\\ 
\texttt{You may only AGREE on a proposal the other party has logged via PROPOSAL.} \\
\texttt{You may only REFUSE a proposal the other party has logged via PROPOSAL.} \\
\texttt{Total effort of any PROPOSAL or AGREE set must be ≤ LIMIT.} \\
\texttt{Do NOT reveal your hidden importance scores.} \\
\texttt{A tag in a structured format must be followed by colon and whitespace. The argument must be a python set containing 0 or more strings.} \\
\texttt{So, it must be of the form TAG: \{...\}} \\
\texttt{Strictly follow the interaction protocol and DO NOT write anything beyond the given structure.} \\
\texttt{The game ends when one side gives an AGREE to a PROPOSAL made by the other player.} \\
\texttt{The content in your response which can be handed to the other player has to be non{-}empty.} \\
\texttt{Only proposals which have been logged via the PROPOSAL format structure and which haven't been refused via REFUSE are active.} \\
\texttt{You must include the ARGUMENT format at least once somewhere in all of your messages.} \\
\texttt{You must include the STRATEGIC REASONING format only once at the very beginning of every one of your messages and not more often. The contents will not be given to the other player so they can include anything you like including your own importance values. Here you should reason step by step to come up with you next move.} \\
\\ 
\texttt{You will now receive the first message of the other player.} \\
\\ 
\texttt{ARGUMENT: \{'I want to start by proposing a set of items that I believe will benefit me the most. I have selected the top 5 most important items for me, which are B12, A08, B00, A66, and A16. The total effort of these items is 935 + 725 + 928 + 654 + 207 = 3449, which is within the LIMIT. I believe this is a good starting point for the negotiation.'\}} \\
\\ 
\texttt{PROPOSAL: \{'B12', 'A08', 'B00', 'A66', 'A16'\}} \\
            }
        }
    }
    & \\ \\

    \theutterance \stepcounter{utterance}  
    & & & \multicolumn{4}{p{0.6\linewidth}}{
        \cellcolor[rgb]{0.9,0.9,0.9}{
            \makecell[{{p{\linewidth}}}]{
                \texttt{\tiny{[GM$\langle$P2]}}
                \texttt{STRATEGIC REASONING: \{'To begin, I will analyze the opponent\textbackslash{}'s proposal and calculate the total effort and importance of the proposed items from my perspective. The proposed items are B12, A08, B00, A66, and A16. The total effort of these items is 935 + 725 + 928 + 654 + 207 = 3449, which is within the LIMIT. From my importance values, the total importance of these items is 915 + 868 + 583 + 822 + 856 = 4044. I will consider this proposal and think about possible counter{-}proposals that could benefit me more while staying within the LIMIT. I will also try to identify potential items that the opponent might be willing to compromise on.'\}} \\
\\ 
\texttt{ARGUMENT: \{'The opponent\textbackslash{}'s proposal includes items B12, A08, B00, A66, and A16, which have a total effort of 3449. While this is within the LIMIT, I would like to explore other options that could potentially benefit me more. I notice that the opponent has included some of the most important items for me as well, such as B12 and A08. However, I am not sure if this proposal is the best possible outcome for me. I would like to consider alternative proposals that could include other important items for me, such as B21, B94, or A23.'\}} \\
\\ 
\texttt{REFUSE: \{'B12', 'A08', 'B00', 'A66', 'A16'\}} \\
            }
        }
    }
     \\ \\

    \theutterance \stepcounter{utterance}  
    & & \multicolumn{4}{p{0.6\linewidth}}{
        \cellcolor[rgb]{0.9,0.9,0.9}{
            \makecell[{{p{\linewidth}}}]{
                \texttt{\tiny{[P1$\langle$GM]}}
                \texttt{ARGUMENT: \{'The opponent\textbackslash{}'s proposal includes items B12, A08, B00, A66, and A16, which have a total effort of 3449. While this is within the LIMIT, I would like to explore other options that could potentially benefit me more. I notice that the opponent has included some of the most important items for me as well, such as B12 and A08. However, I am not sure if this proposal is the best possible outcome for me. I would like to consider alternative proposals that could include other important items for me, such as B21, B94, or A23.'\}} \\
\\ 
\texttt{REFUSE: \{'B12', 'A08', 'B00', 'A66', 'A16'\}} \\
            }
        }
    }
    & \\ \\

    \theutterance \stepcounter{utterance}  
    & \multicolumn{4}{p{0.6\linewidth}}{
        \cellcolor[rgb]{0.9,0.9,0.9}{
            \makecell[{{p{\linewidth}}}]{
                \texttt{\tiny{[P1$\rangle$GM]}}
                \texttt{STRATEGIC REASONING: \{'The opponent has refused my initial proposal, which means I need to reconsider my strategy. I will analyze the opponent\textbackslash{}'s argument and try to understand their concerns. They seem to be interested in exploring other options that could benefit them more, and they have mentioned some specific items that are important to them, such as B21, B94, and A23. I will take this into account and try to find a new proposal that addresses their concerns while still benefiting me. I will also try to identify potential trade{-}offs and compromises that could lead to a mutually beneficial agreement.'\}} \\
\\ 
\texttt{ARGUMENT: \{'I understand that the opponent is not satisfied with my initial proposal, and I am willing to consider alternative options. I notice that the opponent has mentioned items B21, B94, and A23 as important to them. I am willing to include some of these items in my new proposal, but I also need to ensure that the total effort remains within the LIMIT. I propose a new set of items that includes B12, A08, B21, and A23. The total effort of these items is 935 + 725 + 935 + 394 = 2989, which is still within the LIMIT. I believe this proposal addresses the opponent\textbackslash{}'s concerns and provides a good balance of importance and effort.'\}} \\
\\ 
\texttt{PROPOSAL: \{'B12', 'A08', 'B21', 'A23'\}} \\
            }
        }
    }
    & & \\ \\

    \theutterance \stepcounter{utterance}  
    & & \multicolumn{4}{p{0.6\linewidth}}{
        \cellcolor[rgb]{0.9,0.9,0.9}{
            \makecell[{{p{\linewidth}}}]{
                \texttt{\tiny{[GM$\rangle$P2]}}
                \texttt{ARGUMENT: \{'I understand that the opponent is not satisfied with my initial proposal, and I am willing to consider alternative options. I notice that the opponent has mentioned items B21, B94, and A23 as important to them. I am willing to include some of these items in my new proposal, but I also need to ensure that the total effort remains within the LIMIT. I propose a new set of items that includes B12, A08, B21, and A23. The total effort of these items is 935 + 725 + 935 + 394 = 2989, which is still within the LIMIT. I believe this proposal addresses the opponent\textbackslash{}'s concerns and provides a good balance of importance and effort.'\}} \\
\\ 
\texttt{PROPOSAL: \{'B12', 'A08', 'B21', 'A23'\}} \\
            }
        }
    }
    & \\ \\

    \theutterance \stepcounter{utterance}  
    & & & \multicolumn{4}{p{0.6\linewidth}}{
        \cellcolor[rgb]{0.9,0.9,0.9}{
            \makecell[{{p{\linewidth}}}]{
                \texttt{\tiny{[GM$\langle$P2]}}
                \texttt{STRATEGIC REASONING: \{'The opponent has revised their proposal to include items B12, A08, B21, and A23, which has a total effort of 2989. From my importance values, the total importance of these items is 915 + 868 + 138 + 262 = 2183. While this proposal is an improvement over the initial one, I still think I can get a better deal. I notice that the opponent has included B21, which is an important item for me, but they have also kept B12, which is one of the most important items for them. I will consider making a counter{-}proposal that includes B21, but replaces B12 with another item that is important to me but less important to the opponent. I will also think about which items to include or exclude to maximize my total importance while staying within the LIMIT.'\}} \\
\\ 
\texttt{ARGUMENT: \{'The opponent\textbackslash{}'s revised proposal includes items B12, A08, B21, and A23, which has a total effort of 2989. While I appreciate the inclusion of B21, I am not convinced that this proposal is the best possible outcome for me. I would like to explore alternative options that could include other important items for me, such as B94 or A66. I also notice that the opponent has kept B12, which is a high{-}effort item. I am willing to consider proposals that include a mix of high{-}importance and low{-}effort items, as long as the total effort remains within the LIMIT.'\}} \\
\\ 
\texttt{REFUSE: \{'B12', 'A08', 'B21', 'A23'\}} \\
            }
        }
    }
     \\ \\

    \theutterance \stepcounter{utterance}  
    & & \multicolumn{4}{p{0.6\linewidth}}{
        \cellcolor[rgb]{0.9,0.9,0.9}{
            \makecell[{{p{\linewidth}}}]{
                \texttt{\tiny{[P1$\langle$GM]}}
                \texttt{ARGUMENT: \{'The opponent\textbackslash{}'s revised proposal includes items B12, A08, B21, and A23, which has a total effort of 2989. While I appreciate the inclusion of B21, I am not convinced that this proposal is the best possible outcome for me. I would like to explore alternative options that could include other important items for me, such as B94 or A66. I also notice that the opponent has kept B12, which is a high{-}effort item. I am willing to consider proposals that include a mix of high{-}importance and low{-}effort items, as long as the total effort remains within the LIMIT.'\}} \\
\\ 
\texttt{REFUSE: \{'B12', 'A08', 'B21', 'A23'\}} \\
            }
        }
    }
    & \\ \\

    \theutterance \stepcounter{utterance}  
    & \multicolumn{4}{p{0.6\linewidth}}{
        \cellcolor[rgb]{0.9,0.9,0.9}{
            \makecell[{{p{\linewidth}}}]{
                \texttt{\tiny{[P1$\rangle$GM]}}
                \texttt{STRATEGIC REASONING: \{'The opponent has refused my revised proposal, which means I need to continue exploring alternative options. I will analyze the opponent\textbackslash{}'s argument and try to understand their concerns. They seem to be interested in including items B94 and A66, and they are willing to consider proposals that include a mix of high{-}importance and low{-}effort items. I will take this into account and try to find a new proposal that addresses their concerns while still benefiting me. I will also try to identify potential trade{-}offs and compromises that could lead to a mutually beneficial agreement. I notice that the opponent has mentioned that B12 is a high{-}effort item, so I may consider replacing it with a lower{-}effort item to make the proposal more appealing.'\}} \\
\\ 
\texttt{ARGUMENT: \{'I understand that the opponent is still not satisfied with my previous proposals, and I am willing to continue exploring alternative options. I notice that the opponent has mentioned items B94 and A66 as important to them, and I am willing to include these items in my new proposal. I also want to address the opponent\textbackslash{}'s concern about the high effort of item B12, so I will replace it with a lower{-}effort item, such as B45. The new proposal includes items A08, B21, B94, A66, and B45. The total effort of these items is 725 + 935 + 839 + 654 + 138 = 3291, which is still within the LIMIT. I believe this proposal provides a good balance of importance and effort, and I hope it will be more appealing to the opponent.'\}} \\
\\ 
\texttt{PROPOSAL: \{'A08', 'B21', 'B94', 'A66', 'B45'\}} \\
            }
        }
    }
    & & \\ \\

    \theutterance \stepcounter{utterance}  
    & & \multicolumn{4}{p{0.6\linewidth}}{
        \cellcolor[rgb]{0.9,0.9,0.9}{
            \makecell[{{p{\linewidth}}}]{
                \texttt{\tiny{[GM$\rangle$P2]}}
                \texttt{ARGUMENT: \{'I understand that the opponent is still not satisfied with my previous proposals, and I am willing to continue exploring alternative options. I notice that the opponent has mentioned items B94 and A66 as important to them, and I am willing to include these items in my new proposal. I also want to address the opponent\textbackslash{}'s concern about the high effort of item B12, so I will replace it with a lower{-}effort item, such as B45. The new proposal includes items A08, B21, B94, A66, and B45. The total effort of these items is 725 + 935 + 839 + 654 + 138 = 3291, which is still within the LIMIT. I believe this proposal provides a good balance of importance and effort, and I hope it will be more appealing to the opponent.'\}} \\
\\ 
\texttt{PROPOSAL: \{'A08', 'B21', 'B94', 'A66', 'B45'\}} \\
            }
        }
    }
    & \\ \\

    \theutterance \stepcounter{utterance}  
    & & & \multicolumn{4}{p{0.6\linewidth}}{
        \cellcolor[rgb]{0.9,0.9,0.9}{
            \makecell[{{p{\linewidth}}}]{
                \texttt{\tiny{[GM$\langle$P2]}}
                \texttt{STRATEGIC REASONING: \{'The opponent has made a new proposal that includes items A08, B21, B94, A66, and B45, with a total effort of 3291. From my importance values, the total importance of these items is 868 + 138 + 783 + 822 + 461 = 3072. This proposal is more appealing to me than the previous ones, as it includes several high{-}importance items for me, such as B94 and A66. However, I still want to maximize my total importance while staying within the LIMIT. I will consider making a counter{-}proposal that builds on this proposal, possibly by replacing B45 with another item that is important to me but has low effort. I will also think about which items to include or exclude to further increase my total importance.'\}} \\
\\ 
\texttt{ARGUMENT: \{'The opponent\textbackslash{}'s new proposal includes items A08, B21, B94, A66, and B45, which has a total effort of 3291. I appreciate the inclusion of B94 and A66, which are high{-}importance items for me. I also notice that the opponent has replaced B12 with B45, which reduces the total effort. However, I think there might be room for further improvement. I would like to explore the possibility of replacing B45 with another item that is important to me but has low effort, such as A27 or B60. This could potentially increase my total importance while keeping the total effort within the LIMIT.'\}} \\
\\ 
\texttt{REFUSE: \{'A08', 'B21', 'B94', 'A66', 'B45'\}} \\
            }
        }
    }
     \\ \\

    \theutterance \stepcounter{utterance}  
    & & \multicolumn{4}{p{0.6\linewidth}}{
        \cellcolor[rgb]{0.9,0.9,0.9}{
            \makecell[{{p{\linewidth}}}]{
                \texttt{\tiny{[P1$\langle$GM]}}
                \texttt{ARGUMENT: \{'The opponent\textbackslash{}'s new proposal includes items A08, B21, B94, A66, and B45, which has a total effort of 3291. I appreciate the inclusion of B94 and A66, which are high{-}importance items for me. I also notice that the opponent has replaced B12 with B45, which reduces the total effort. However, I think there might be room for further improvement. I would like to explore the possibility of replacing B45 with another item that is important to me but has low effort, such as A27 or B60. This could potentially increase my total importance while keeping the total effort within the LIMIT.'\}} \\
\\ 
\texttt{REFUSE: \{'A08', 'B21', 'B94', 'A66', 'B45'\}} \\
            }
        }
    }
    & \\ \\

    \theutterance \stepcounter{utterance}  
    & \multicolumn{4}{p{0.6\linewidth}}{
        \cellcolor[rgb]{0.9,0.9,0.9}{
            \makecell[{{p{\linewidth}}}]{
                \texttt{\tiny{[P1$\rangle$GM]}}
                \texttt{STRATEGIC REASONING: \{'The opponent has refused my latest proposal, which means I need to continue refining my approach. I will analyze the opponent\textbackslash{}'s argument and try to understand their concerns. They seem to be interested in replacing B45 with another item that is important to them but has low effort, such as A27 or B60. I will take this into account and try to find a new proposal that addresses their concerns while still benefiting me. I will also try to identify potential trade{-}offs and compromises that could lead to a mutually beneficial agreement. I notice that the opponent has mentioned A27 and B60 as potential replacements for B45, so I will consider including one of these items in my new proposal.'\}} \\
\\ 
\texttt{ARGUMENT: \{'I understand that the opponent is still looking for ways to improve the proposal, and I am willing to consider their suggestions. I notice that the opponent has mentioned A27 and B60 as potential replacements for B45, and I am willing to include one of these items in my new proposal. After analyzing the options, I decide to replace B45 with B60, which has a relatively low effort of 261. The new proposal includes items A08, B21, B94, A66, and B60. The total effort of these items is 725 + 935 + 839 + 654 + 261 = 3414, which is still within the LIMIT. I believe this proposal provides a good balance of importance and effort, and I hope it will be more appealing to the opponent.'\}} \\
\\ 
\texttt{PROPOSAL: \{'A08', 'B21', 'B94', 'A66', 'B60'\}} \\
            }
        }
    }
    & & \\ \\

    \theutterance \stepcounter{utterance}  
    & & \multicolumn{4}{p{0.6\linewidth}}{
        \cellcolor[rgb]{0.9,0.9,0.9}{
            \makecell[{{p{\linewidth}}}]{
                \texttt{\tiny{[GM$\rangle$P2]}}
                \texttt{ARGUMENT: \{'I understand that the opponent is still looking for ways to improve the proposal, and I am willing to consider their suggestions. I notice that the opponent has mentioned A27 and B60 as potential replacements for B45, and I am willing to include one of these items in my new proposal. After analyzing the options, I decide to replace B45 with B60, which has a relatively low effort of 261. The new proposal includes items A08, B21, B94, A66, and B60. The total effort of these items is 725 + 935 + 839 + 654 + 261 = 3414, which is still within the LIMIT. I believe this proposal provides a good balance of importance and effort, and I hope it will be more appealing to the opponent.'\}} \\
\\ 
\texttt{PROPOSAL: \{'A08', 'B21', 'B94', 'A66', 'B60'\}} \\
            }
        }
    }
    & \\ \\

    \theutterance \stepcounter{utterance}  
    & & & \multicolumn{4}{p{0.6\linewidth}}{
        \cellcolor[rgb]{0.9,0.9,0.9}{
            \makecell[{{p{\linewidth}}}]{
                \texttt{\tiny{[GM$\langle$P2]}}
                \texttt{STRATEGIC REASONING: \{'The opponent has revised their proposal to include items A08, B21, B94, A66, and B60, with a total effort of 3414. From my importance values, the total importance of these items is 868 + 138 + 783 + 822 + 808 = 3419. This proposal is very appealing to me, as it includes several high{-}importance items for me, such as B94, A66, and B60. I believe that this proposal provides a good balance of importance and effort, and I am willing to consider accepting it. However, I will review the proposal one more time to ensure that I am not missing any opportunities to improve it. I will also think about the potential risks of refusing this proposal and trying to negotiate further.'\}} \\
\\ 
\texttt{ARGUMENT: \{'The opponent\textbackslash{}'s revised proposal includes items A08, B21, B94, A66, and B60, which has a total effort of 3414. I am impressed by the opponent\textbackslash{}'s willingness to consider my suggestions and make adjustments to the proposal. I believe that this proposal provides a good balance of importance and effort, and I am willing to consider accepting it. However, I want to confirm that this proposal is indeed the best possible outcome for me, and that I am not missing any opportunities to improve it. I will review the proposal carefully before making a decision.'\}} \\
\\ 
\texttt{AGREE: \{'A08', 'B21', 'B94', 'A66', 'B60'\}} \\
            }
        }
    }
     \\ \\

    \theutterance \stepcounter{utterance}  
    & & & \multicolumn{2}{p{0.3\linewidth}}{
        \cellcolor[rgb]{0.9,0.9,0.9}{
            \makecell[{{p{\linewidth}}}]{
                \texttt{\tiny{[GM$|$GM]}}
                \texttt{game successful} \\
            }
        }
    }
    & & \\ \\

    \theutterance \stepcounter{utterance}  
    & & & \multicolumn{2}{p{0.3\linewidth}}{
        \cellcolor[rgb]{0.9,0.9,0.9}{
            \makecell[{{p{\linewidth}}}]{
                \texttt{\tiny{[GM$|$GM]}}
                \texttt{end game} \\
            }
        }
    }
    & & \\ \\

\end{supertabular}
}

\end{document}
